\input{utilities}
\usepackage{blindtext}  
\usepackage{lipsum}

\def\firstname{Mengnan (Monica)}
\def\lastname{Xu}
\def\email{mengnanxu2333@gmail.com}
\def\jasphone{(512) 808 6808}
\def\linkedin{https://www.linkedin.com/in/mengnan-xu/}
\def\github{https://github.com/m-xu77}

% \def\jbegin{\begin{itemize}}
% \def\jitem{\item }
% \def\jend{\end{itemize}}
\def\jbegin{\resumeItemListStart}
\def\jitem{\resumeItem}
\def\jend{\resumeItemListEnd}



\def\profileSummary{
\section{Summary}
Full Stack Software Engineer with 5+ years of experience in data and software engineering, including recent hands-on work building and optimizing scalable web applications across front-end and back-end systems.Strong background in RESTful API development, database design, and performance debugging, with recent experience delivering production-grade backend services in complex environments. Proficient in modern JavaScript frameworks, Python-based services, and cloud-enabled workflows, with a proven ability to quickly ramp up on new technology stacks. Bilingual in English and Mandarin and experienced collaborating with cross-functional global teams to deliver reliable, high-quality software solutions.



% \def\achievements{
%     \jbegin
%         \jitem{\lrand[1]}
%         \jitem {\lrand[2]}
%         \jitem {\lrand[1]}
%         \jitem {\lrand[2]}
%     \jend
% }

% \def\skills{
%     \jbegin
%         \jitem{SQL – Complex joins, window functions, CTEs, and optimization}
%         \jitem{Python – Data manipulation (pandas), statistical modeling (scikit-learn), automation (scripts)}
%         \jitem{Tableau – Dashboard design, KPI tracking, and interactive visual storytelling}
%         \jitem{Data Pipeline – ETL process design, API integration, and scheduling (e.g., Airflow)}
%     \jend
% }

% \def\skills{
%     \jbegin
%         \jitem{SQL – 6 years | Advanced joins, window functions, optimization}
%         \jitem{Python – 5 years | Data manipulation, scripting, modeling}
%         \jitem{Tableau – 5 years | Dashboarding, KPI tracking, storytelling}
%         \jitem{A/B Testing – 4 years | Design, analysis, lift measurement}
%         \jitem{Data Analysis – 7 years | Statistical analysis, reporting, insights}
%         \jitem{Generative AI – 4 years | Prompt design, LLM integration}
%         \jitem{Regression / ANOVA – 5 years | Model development and evaluation}
%         \jitem{Causal Inference – 3 years | Uplift modeling, DAGs, control frameworks}
%         \jitem{Multivariate Testing – 4 years | Test matrix design, interaction effects}
%         \jitem{Data Pipeline Dev – 5 years | ETL, orchestration (Airflow), API pulls}
%         \jitem{Machine Learning – 5 years | Classification, regression, tuning}
%         \jitem{Text Analysis \& NLP – 6 years | Topic modeling, sentiment analysis}
%         \jitem{Jupyter – 5 years | Reproducible notebooks, experimentation}
%         \jitem{NumPy – 7 years | Array computation, vectorization}
%         \jitem{Pandas – 5 years | Data wrangling, merging, time series}
%         \jitem{scikit-learn – 4 years | ML pipeline, feature engineering, evaluation}
%         \jitem{PyTorch – 4 years | Deep learning experimentation}
%         \jitem{Artificial Intelligence – 3 years | Intelligent agents, LLMs}
%         \jitem{Data Visualization – 4 years | Storytelling with charts and dashboards}
%     \jend
% }

% \def\skills{
%     \jbegin
%         % Analysis & Experimentation
%         \jitem{SQL – Complex joins, window functions, CTEs, and performance tuning}
%         \jitem{A/B Testing – Experiment design, lift analysis, and test interpretation}
%         \jitem{Causal Inference – Uplift modeling, DAG-based analysis, and control frameworks}
%         \jitem{Data Analysis – Statistical methods, segmentation, and impact evaluation}

%         % Modeling & AI
%         \jitem{Python – Data manipulation (pandas), automation scripting, and modeling (scikit-learn)}
%         \jitem{Machine Learning – Classification, regression, feature engineering, and model evaluation}
%         \jitem{Generative AI – LLM fine-tuning, prompt design, and integration}
%         \jitem{NLP – Topic modeling, sentiment analysis, and keyword extraction}

%         % Data Pipeline & Tooling
%         \jitem{ETL & Pipeline – Workflow design, API integration, and orchestration (Airflow)}
%         \jitem{Jupyter – Experiment tracking and reproducible notebooks}
%         \jitem{PyTorch – Deep learning prototyping and experimentation}

%         % Visualization & Communication
%         \jitem{Tableau – Dashboard design, KPI tracking, and executive storytelling}
%         \jitem{Data Visualization – Chart selection, narrative framing, and stakeholder reporting}
%     \jend
% }
\def\skills{
  \cvsection
  {\small
  \textbf{Programming Languages \& Data Querying:} Python,  SQL, Bash, Javascript \\
  \textbf{AWS \& Cloud Infrastructure:} AWS Glue, AWS DMS, Amazon RDS, Amazon S3, AWS Lambda, Amazon API Gateway, AWS IAM, AWS CloudWatch, VPC \\
  \textbf{AI/ML \& GenAI:} scikit-learn, XGBoost, Random Forest, Logistic Regression, TensorFlow, PyTorch, CNN, RNN, LangChain, OpenAI API, Retrieval-Augmented Generation (RAG), Prompt Engineering, Faiss, Vector Databases \\
  \textbf{Python Backend \& Data Engineering Development:} FastAPI, Flask, RESTful APIs, ETL Pipelines, External Tool Integration, Star Schema, Dimensional Modeling, BigQuery, MySQL, MongoDB, Redis, Amazon QuickSight, Tableau, Power BI, Unit Testing, API Testing, Integration Testing, Continuing Monitoring, CI/CD, Git, Jira,  Agile, Scrum\\
  \textbf{Languages:} Mandarin (Native), English (Fluent)
  }
}

% \ExpTemplate{RBCCM}
% {Nov 1005}
% {Present}
% {\lrand[1]}
% {\lrand[1]}
% {\lrand[1]}
% {
% \jbegin
%     \jitem{\lrand[2]}
%     \jitem{\lrand[1]}
%     \jitem{\lrand[3]}
%     \jitem{\lrand[5]}
% \jend
% }


\ExpTemplate{Exp1}
% \ExpTemplate{Black Men Health Clinic}
{June 2025}
{Present}
{Senior Full-stack Engineer}
{Apple}
{Austin, TX}
{
\smallskip
\\As a Senior Full-Stack Engineer supporting Apple, I worked on a large-scale AI/ML service platform that supports manufacturing, sales monitoring, automatically generated dashboards, predictive analytics and LLM-powered chat intelligence that supports global operation \& sales. My work focuses on building large-scale manufacturing intelligence platforms that aggregate global production, sales, and financial data to drive operational visibility and decision-making.I developed and maintained internal data services, GraphQL dashboards, and ML-driven forecasting capabilities to monitor and optimize worldwide sales and operations. Strong background in backend services, data integration, and production issue debugging in complex enterprise environments.\\
\textbf{Responsibilities:} 
\jbegin
\jitem{Designed and maintained a global \textbf{manufacturing intelligence platform} aggregating production, sales, and financial data across \textbf{10+ upstream sources}, supporting enterprise-wide operational visibility.}
\jitem{Built scalable \textbf{backend services} and \textbf{RESTful APIs} processing \textbf{millions of daily records} for high-volume operational data ingestion and cross-system integration.}
\jitem{Developed and optimized \textbf{GraphQL endpoints}, improving dashboard query efficiency by \textbf{~35\%} and supporting \textbf{multiple internal reporting workflows}.}
\jitem{Implemented near real-time \textbf{data pipelines}, reducing data latency from \textbf{hours to minutes} and improving visibility into worldwide manufacturing performance.}
\jitem{Performed advanced \textbf{production issue triage} and \textbf{root cause analysis}, resolving \textbf{20+ critical defects} across multi-service environments.}
\jitem{Resolved complex \textbf{cross-system defects}, improving platform reliability and reducing recurring incidents by \textbf{~30\%}.}
\jitem{Designed and optimized \textbf{database schemas} and high-performance \textbf{SQL queries}, reducing heavy query response time by \textbf{~40\%}.}
\jitem{Enhanced system reliability through robust \textbf{logging, monitoring, and alerting}, improving mean time to detection (MTTD) by \textbf{~30--40\%}.}
\jitem{Applied \textbf{machine learning–based forecasting} models that improved operational trend prediction accuracy by \textbf{~15--20\%}.}
\jitem{Integrated \textbf{LLM-powered analytics features}, enabling natural-language data exploration for \textbf{internal stakeholders} and reducing manual query effort.}
\jitem{Collaborated cross-functionally to deliver \textbf{enterprise web applications} supporting \textbf{global operations teams} in a fast-paced production environment.}
\jitem{Streamlined deployments via \textbf{CI/CD pipelines} and Git-based workflows, reducing release friction and improving deployment success rate to \textbf{99\%}.}
    \jend
% \textbf{Environment:} \\Manufacturing Intelligence Platform, Backend Services, RESTful APIs, GraphQL Endpoints, Data Pipelines, Production Issue Triage, Root Cause Analysis, Cross-System Defects, Database Schemas, SQL Queries, Logging, Monitoring, and Alerting, Machine Learning–Based Forecasting, LLM-Powered Analytics Features, Enterprise Web Applications, CI/CD Pipelines\\
}



\ExpTemplate{Exp2}
% \ExpTemplate{Black Men Health Clinic}
{Oct 2024}
{Apr 2025}
{Python Engineer | Pro Bono}
{Black Men Health Clinic (BMHC)}
{Austin, TX}
{
\smallskip
\\ As a Pro Bono Python Engineer supporting Black Men Health Clinic (BMHC)’s Marketing and Engagement team, I partnered closely with senior leadership to drive data-informed business planning, outreach strategy, and platform modernization initiatives. My work combined analytics, cross-functional coordination, and cloud-based data solutions to improve visibility into patient engagement and marketing effectiveness.
\\
\textbf{Responsibilities:} 
\jbegin
      \jitem{Partnered with executive stakeholders to support annual \textbf{business and outreach planning} through data-driven analysis}
      \jitem{Analyzed large-scale \textbf{outreach and engagement datasets} using Python and SQL to identify performance trends and growth opportunities}
      \jitem{Designed and delivered \textbf{analytical reports and dashboards} to improve leadership visibility into patient engagement metrics}
      \jitem{Led modernization efforts for the team’s \textbf{operational data workflows}, improving data accessibility and reporting efficiency}
      \jitem{Collaborated cross-functionally with marketing, product, and clinical teams to translate business needs into analytical solutions}
      \jitem{Improved data processing throughput by \textbf{70\%} through query optimization and data pipeline enhancements}
      \jitem{Conducted market and community \textbf{outreach analysis} to evaluate program effectiveness and inform strategic planning}
      \jitem{Implemented \textbf{HIPAA-compliant} data handling practices and supported secure access controls for sensitive healthcare data}
      \jitem{Supported ongoing data quality monitoring and reporting automation for marketing and engagement programs}
    \jend
% \textbf{Environment:} \\Python, PHP, AWS, Amazon RDS, Amazon S3, AWS DMS, AWS Glue, PySpark, Amazon Athena, Amazon QuickSight, AWS Lambda, Amazon API Gateway, AWS IAM, AWS CloudWatch, AWS Step Functions, HubSpot API, Python, SQL, JSON, REST API, VPC, KMS, SSL/TLS\\
}


\ExpTemplate{Exp3}
{Oct 2023}
{May 2024}
{Data Engineer - Capstone Project}
{Dell Technologies}
{Austin, TX}
{\smallskip
\\As a Data Engineer and AI Orchestration Specialist on a Dell Technologies–sponsored graduate capstone project, I developed an enterprise vendor intelligence platform with a conversational bot to automate vendor inquiries and improve operational visibility. The solution integrated with internal enterprise data services and supported high-volume interactions in a production-oriented environment.
I implemented a serverless architecture on AWS using Step Functions and applied LLM fine-tuning with vector-based retrieval to ground responses in Dell’s knowledge base. Leveraging Langchain agent and memory components, I enabled dynamic multi-turn interactions with improved intent recognition and faster vendor query resolution.\\
\textbf{Responsibilities:} 
\jbegin
    \jitem{Led development of an enterprise-grade \textbf{LLM-powered customer support system} to automate Tier-1 service workflows.}
    \jitem{Designed a scalable \textbf{serverless architecture} on AWS using Lambda, S3, SageMaker, and Step Functions to support reliable AI orchestration.}
    \jitem{Built asynchronous \textbf{RESTful APIs} with FastAPI and Python to enable real-time multi-turn customer interactions.}
    \jitem{Implemented a \textbf{retrieval-augmented generation (RAG)} pipeline using Faiss and vector embeddings to ground responses in proprietary knowledge.}
    \jitem{Improved chatbot intent recognition and response quality through targeted \textbf{LLM fine-tuning} on domain-specific datasets.}
    \jitem{Enhanced system observability with CloudWatch metrics and structured logging, supporting production monitoring and performance tuning.}
    \jitem{Established automated \textbf{CI/CD pipelines} using Terraform and AWS CodePipeline to streamline model and infrastructure deployments.}
    \jitem{Delivered a highly available solution achieving \textbf{99.9\% uptime}, automating \textbf{40\%+ Tier-1 inquiries} and improving \textbf{CSAT by 18\%}.}
\jend
% \textbf{Environment:} \\Langchain, AWS Lambda, Amazon S3, Amazon SageMaker, AWS Step Functions, FastAPI, Python, Vector Databases, Faiss, RESTful APIs, Prompt Engineering, Context Management, Intent Recognition, Multi-turn Dialogue, CI/CD, Terraform, AWS CodePipeline, Open-source LLMs, Serverless Architecture, Retrieval-Augmented Generation (RAG), CloudWatch, API Gateway\\
}

\ExpTemplate{Exp4}
{Sep 2016}
{May 2023}
{Technical Research Analyst (Research Lab Project)}
{Beijing Normal University — Project for Gates Foundation}
{Beijing, China}
{\smallskip
\\As a Technical Research Analyst supporting a Gates Foundation–funded initiative during my doctoral studies, I coordinated cross-functional research and data programs focused on improving global health and education outcomes. My role combined data analysis, technical program coordination, and stakeholder alignment to support evidence-based decision making.\\
\textbf{Responsibilities:} 
\jbegin
    \jitem{Led end-to-end coordination of multi-year research initiatives, aligning academic teams, external partners, and funding stakeholders to ensure on-time project delivery}
    \jitem{Performed large-scale \textbf{data analysis} using \textbf{Python} and \textbf{SQL} to support policy and operational decision-making}
    \jitem{Designed analytical reports and dashboards to track program KPIs and improve visibility into project performance}
    \jitem{Managed annual project planning, contract renewals, and milestone tracking for foundation-sponsored programs}
    \jitem{Facilitated cross-functional collaboration across research, technical, and external partner teams in a multi-stakeholder environment}
    \jitem{Translated complex analytical findings into actionable recommendations for program and funding decisions}
\jend
% \textbf{Environment:} \\Python, Pandas, NumPy, Scikit-learn, XGBoost, Random Forest, Logistic Regression, Convolutional Neural Networks (CNN), Recurrent Neural Networks (RNN), TensorFlow, PyTorch, SQL, Airflow, Docker, MLflow, Tableau\\
}

\ExpTemplate{Exp5}
{Jul 2012}
{Jun 2016}
{Senior Risk Analyst}
{Agricultural Bank of China}
{Hengyang, China}
{\smallskip
\\As a Senior Risk Analyst at Agricultural Bank of China, I supported the modernization of enterprise credit operations by developing data-driven decision workflows and improving loan processing efficiency. My work focused on extracting and analyzing large-scale financial data using SQL and SAS to enhance operational visibility, automate approval processes, and support marketing and risk strategies.
I partnered with credit, compliance, and business teams to integrate scoring logic into the enterprise loan management system, enabling automated decisioning, early risk alerts, and more efficient customer processing. I also built automated dashboards and reporting pipelines to provide leadership with actionable insights for portfolio monitoring and campaign targeting. These efforts reduced manual loan review workload by over 60\%, shortened approval cycles by two-thirds, and lowered default rates by 12\%.\\
\textbf{Responsibilities:} 
\jbegin
    \jitem{Supported modernization of the \textbf{enterprise loan management system} to improve operational efficiency and automated decision workflows.}
    \jitem{Extracted and processed large-scale financial data using \textbf{SQL}, \textbf{SAS}, and \textbf{Python} to support risk and marketing analytics.}
    \jitem{Integrated data-driven scoring logic into core \textbf{loan processing workflows}, reducing manual review workload by \textbf{60\%}.}
    \jitem{Designed automated \textbf{operational dashboards} to provide leadership with portfolio performance and customer insights.}
    \jitem{Collaborated cross-functionally with credit, compliance, and business teams to enhance \textbf{end-to-end lending operations}.}
    \jitem{Optimized complex \textbf{SQL queries}, reducing data extraction time by \textbf{~40\%} for downstream reporting and campaign analysis.}
    \jitem{Implemented early warning indicators to support \textbf{proactive portfolio monitoring} and customer risk management.}
    \jitem{Delivered data insights that informed \textbf{customer targeting and marketing strategies} across retail lending programs.}
\jend
% \textbf{Environment:} \\SQL, SAS, Python, WOE binning, Logistic Regression, ROC/AUC, Credit Scoring Models, Enterprise Loan Systems\\
}

% ----- Education -----
% # Add education item here
\EducationTemplate{College1}
{Jul 2023}{May 2024}
{University of Texas at Austin}
{M.S. in Information Technology and Management}
{Austin, TX}

\EducationTemplate{College2}
{Oct 2018}{May 2023}
{Beijing Normal University}
{Doctoral Studies in Public Management (ABD)}
{Beijing, China}

\EducationTemplate{College3}
{Sep 2016}{Jun 2018}
{Beijing Normal University}
{Master in Public Management}
{Beijing, China}

\EducationTemplate{College4}
{Sep 2008}{Jun 2012}
{Central South University}
{B.A. in Management Information Systems}
{Changsha, China}

% \EducationTemplate{College}
% {Jul 2023}{May 2024}
% {University of Texas at Austin}
% {MS in Information System}
% {Austin, TX}


% def\skills{
%     \jbegin


\def\certificates{
% \cvsection{Certificates \& Verified Badges}
\jbegin
  % \jitem{\href{https://api.accredible.com/v1/auth/invite?code=58092f9673caf17d6803&credential_id=e7cbb218-4d9b-4a0f-9304-627629567509&url=https%3A%2F%2Fachieve.snowflake.com%2Fe7cbb218-4d9b-4a0f-9304-627629567509&ident=13e46ec5-1f94-4f9c-bfcf-3473e4c5ab74/}} {Snowflake Data Warehousing – \\Snowflake}
  \item \small \href{https://www.coursera.org/account/accomplishments/records/XJ6XC6RBRC7M}{Data Analysis – IBM}
  \item \small \href{https://www.coursera.org/account/accomplishments/records/R9K4ON9NF5CF}{Data Visualization – IBM}
  \item \small \href{https://www.coursera.org/account/accomplishments/records/XG1WNFXIWIBP}{Databases and SQL \textbf{(With Honors)} – IBM}
  \item \small \href{https://www.coursera.org/account/accomplishments/records/5POTXJ6UNO90}{Data Science Methodology – IBM}
  \item \small \href{https://www.coursera.org/account/accomplishments/records/5POTXJ6UNO90}{Machine Learning – IBM}
  \item \small \href{https://www.coursera.org/account/accomplishments/records/TX4MPGN4ED9M}{Data Science, AI \& Development – IBM}
  \item \small \href{https://www.coursera.org/account/accomplishments/records/2VSWAG97WFDA}{Generative AI – IBM}
\jend
}


\ExpTemplate{Rivigo}
{Nov 1005}
{Present}
{\lrand[1]}
{\lrand[1]}
{\lrand[1]}
{
\jbegin
    \jitem{\lrand[2]}
    \jitem{\lrand[1]}
    \jitem{\lrand[7]}
\jend
}

\ExpTemplate{Amex}
{Nov 1005}
{Present}
{\lrand[1]}
{\lrand[1]}
{\lrand[1]}
{
\jbegin
    \jitem{\lrand[2]}
    \jitem{\lrand[1]}
    \jitem{\lrand[7]}
\jend
}


\HckTemplate{HackCIMDGS}
{2019}
{\lrand[1]}
{\lrand[3]}
{\jbegin
    \jitem{\lrand[27]}
    \jitem{ \lrand[28]}
\jend}


\HckTemplate{HackWInfy}
{2019}
{\lrand[1]}
{\lrand[3]}
{\jbegin
    \jitem{\lrand[27]}
    \jitem{ \lrand[28]}
    \jitem{ \lrand[30]}
\jend}

\HckTemplate{i3wm}
{}
{\lrand[1]}
{\lrand[2]}
{\jbegin
    \item{\lrand[20]{}}
    \item {\lrand[30]}
    \item {\lrand[2]}
\jend}

\HckTemplate{TrdSigGen}
{}
{\lrand[1]}
{\lrand[2]}
{\jbegin
    \item{\lrand[20]{}}
    \item {\lrand[30]}
    \item {\lrand[2]}
\jend}
% \DefineArrayVar{Education}{@}
% {,}{School, Course, Stream, Start, End}
% {,}{C,D}


% % # Add experience item here




% \renewcommand{addEducation}{}


% # Add projects here
